\section{Conclusiones y límites}\label{sec:conclusiones}

\subsection{Conclusiones}

Con los parámetros ensayados, y siempre como comparación entre escenarios:

\begin{itemize}
  \item \textbf{Entrada (P1, P2).} Los tornos son el recurso más cargado, con una
    intensidad de tráfico de 0,93 en la punta. Tres tornos por puerta, el valor
    actual, cumplen el objetivo de espera y dos no; el margen es sensible al
    tiempo de servicio.
  \item \textbf{Llegada anticipada (P3).} Reduce la espera y, sobre todo, la
    entrada tardía, que depende más de cuándo llega el público y del recorrido
    hasta el asiento que de la cola del torno.
  \item \textbf{Prioridad (P4).} Beneficia a mayores y PMR con un coste pequeño y
    repartido para el resto, sin cambiar la espera total.
  \item \textbf{Salida escalonada (P5).} Reduce casi a la mitad la aglomeración en
    las salidas a cambio de alargar algo más de un minuto el vaciado.
  \item \textbf{Acomodadores (P6).} Son decisivos cuando se duplica la probabilidad
    de conflicto: con 3 por macrozona el sistema se satura y con 9 la duración de
    un conflicto se acerca a la de su resolución.
  \item \textbf{Cierre de puerta (P7).} La redirección a la puerta con menos cola
    absorbe casi todo el efecto de un cierre en la punta, en condiciones
    favorables a la redirección.
  \item \textbf{Evacuación (P8).} Con todas las salidas la capacidad basta, pero no
    siempre se evacúa a tiempo a las personas con movilidad reducida; con una
    salida bloqueada no se cumple la referencia de 8~minutos. El promedio de
    vaciado, por sí solo, no responde a la pregunta.
\end{itemize}

Estas conclusiones deben acompañarse de los IC, los tamaños por escenario y los
límites siguientes.

\subsection{Límites}

\begin{itemize}
  \item \textbf{Validación externa.} La verificación comprueba conservación,
    orden, capacidades, casos límite, cálculo de indicadores y reproducibilidad.
    No hay datos completos de colas, salida y evacuación con los que cuantificar
    el error respecto al estadio real. El ritmo por torno se ha contrastado con la
    cifra que la Parte 1 atribuye a la fuente F2, sin volver a comprobarla en el
    artículo original.
  \item \textbf{Precisión.} El criterio de precisión controla la incertidumbre
    muestral de la espera media; no controla el error de los parámetros, de la
    representación espacial ni de las reglas de comportamiento.
  \item \textbf{Generalización.} Las decisiones se estudian con perfiles,
    probabilidades y recursos concretos: no cabe generalizar que tres tornos
    basten para cualquier demanda ni que nueve acomodadores sean siempre
    suficientes.
  \item \textbf{Representación espacial.} El exterior del estadio no está
    representado y la redirección no añade marcha ni información imperfecta.
    Los peatones tienen velocidades y distancias sorteadas, sin un modelo de
    densidad, empujones o asistencia. Esto afecta sobre todo a P7 y P8.
  \item \textbf{Hipótesis adoptadas.} En particular, que el cierre de una puerta
    solo afecta a quien llega mientras dura (sección~\ref{sec:desviaciones}).
    No influye con el cierre estudiado, pero no se ha comprobado con cierres
    posteriores ni con una redirección menos favorable.
  \item \textbf{Sensibilidad.} S1 cambia cada factor por separado y no incluye la
    emergencia; no responde a combinaciones de aforo, servicio y velocidad ni a
    medidas de evacuación asistida.
\end{itemize}
